% GENERATED FILE. Edit canonical lessons, then run npm run generate:books.
% canonical-source-hash: fnv1a64:b6fb26505e780244
% canonical-lessons: HI-C447-kunj, HI-C447-bel, HI-C447-kai, HI-C447-kira, HI-C447-lipatna

\chapter{Bushes and Creepers}
\label{ch:hi-bushes-and-creepers}
\hlchaptermodality{447}

\begin{chapteropening}
\textbf{By the end of this chapter:} \emph{I can talk about bushes, creepers, moss and insects, and a vine winding round a tree.}

Complete the last lesson of chapter 447: I can talk about bushes, creepers, moss and insects, and a vine winding round a tree.
\end{chapteropening}

\section[\texorpdfstring{\hi{कुंज} (kuñj)}{कुंज (kuñj)}]{\hi{कुंज} (kuñj) — a bower, a grove}

\label{lesson:HI-C447-kunj}



\begin{quote}
Before the new one: say the Hindi for to bloom, then the Hindi for to wilt.
\end{quote}

\subsection*{You'll want to know: \hi{कुंज}}
\textbf{\hi{कुंज}} — \emph{kuñj} — \textquotedblleft{}a bower, a grove\textquotedblright{}.

\begin{quote}
\hi{बगीचे} \hi{के} \hi{कोने} \hi{में} \hi{एक} \hi{छोटा} \hi{कुंज} \hi{है।}
\end{quote}

\begin{quote}
\emph{bagīce ke kone meṁ ek choṭā kuñj hai.}
\end{quote}

\begin{quote}
There is a little bower in the corner of the garden.
\end{quote}

Masculine.

\begin{cousinweb}
\textbf{\hi{कुंज}} (\emph{kuñj}) is a tatsama, Sanskrit \emph{kuñja}, \textquotedblleft{}a place overrun with plants or creepers, a bower\textquotedblright{} (Monier-Williams).
\end{cousinweb}

\begin{grammarlens}[title={What you've built}]
One more word for getting around.
\end{grammarlens}

\subsection*{Your turn}
\begin{itemize}
\raggedright
  \item \emph{Say it:} \emph{kuñj}
  \item \emph{Say it:} \emph{kuñj}, once more
  \item \emph{Say it:} \emph{bagīce ke kone meṁ ek choṭā kuñj hai.} — There is a little bower in the corner of the garden.
  \item \emph{Recall:} say the Hindi for to bloom, then the Hindi for to wilt, then say \emph{kuñj} again
\end{itemize}

\begin{tcolorbox}[breakable,colback=teal!4,colframe=teal!35!black,title={Before you move on}]
What does \hi{कुंज} mean? (\textquotedblleft{}A bower, a grove\textquotedblright{}.) Say it once more.
\end{tcolorbox}
\section[\texorpdfstring{\hi{बेल} (bel)}{बेल (bel)}]{\hi{बेल} (bel) — a creeper, a vine}

\label{lesson:HI-C447-bel}



\begin{quote}
Before the new one: say the Hindi for to wilt, then the Hindi for a bower.
\end{quote}

\subsection*{You'll want to know: \hi{बेल}}
\textbf{\hi{बेल}} — \emph{bel} — \textquotedblleft{}a creeper, a vine\textquotedblright{}.

\begin{quote}
\hi{दीवार} \hi{पर} \hi{हरी} \hi{बेल} \hi{चढ़} \hi{रही} \hi{है।}
\end{quote}

\begin{quote}
\emph{dīvār par harī bel caṛh rahī hai.}
\end{quote}

\begin{quote}
A green creeper is climbing the wall.
\end{quote}

Feminine.

A different \hi{बेल}, masculine, is the wood-apple tree.

\begin{cousinweb}
\textbf{\hi{बेल}} (\emph{bel}) is a tadbhava of Sanskrit \emph{velli}, \textquotedblleft{}a creeping plant\textquotedblright{} (CDIAL 12123), which lists Hindi \emph{bel}, \textquotedblleft{}creeper\textquotedblright{}.
\end{cousinweb}

\begin{grammarlens}[title={What you've built}]
One more word for getting around.
\end{grammarlens}

\subsection*{Your turn}
\begin{itemize}
\raggedright
  \item \emph{Say it:} \emph{bel}
  \item \emph{Say it:} \emph{bel}, once more
  \item \emph{Say it:} \emph{dīvār par harī bel caṛh rahī hai.} — A green creeper is climbing the wall.
  \item \emph{Recall:} say the Hindi for to wilt, then the Hindi for a bower, then say \emph{bel} again
\end{itemize}

\begin{tcolorbox}[breakable,colback=teal!4,colframe=teal!35!black,title={Before you move on}]
What does \hi{बेल} mean? (\textquotedblleft{}A creeper, a vine\textquotedblright{}.) Say it once more.
\end{tcolorbox}
\section[\texorpdfstring{\hi{काई} (kāī)}{काई (kāī)}]{\hi{काई} (kāī) — moss; green scum (on water)}

\label{lesson:HI-C447-kai}



\begin{quote}
Before the new one: say the Hindi for a bower, then the Hindi for a creeper.
\end{quote}

\subsection*{You'll want to know: \hi{काई}}
\textbf{\hi{काई}} — \emph{kāī} — \textquotedblleft{}moss; green scum (on water)\textquotedblright{}.

\begin{quote}
\hi{पत्थरों} \hi{पर} \hi{काई} \hi{थी}, \hi{इसलिए} \hi{मैं} \hi{गिर} \hi{गया।}
\end{quote}

\begin{quote}
\emph{pattharoṁ par kāī thī, isliye maiṁ gir gayā.}
\end{quote}

\begin{quote}
There was moss on the stones, so I fell over.
\end{quote}

Feminine.

\begin{cousinweb}
\textbf{\hi{काई}} (\emph{kāī}) goes back to \emph{*kāvikā}, \textquotedblleft{}scum\textquotedblright{}, a form CDIAL reconstructs (CDIAL 3109), which lists Hindi \emph{kāī}, \textquotedblleft{}green scum on stagnant water\textquotedblright{}.
\end{cousinweb}

\begin{grammarlens}[title={What you've built}]
One more word for getting around.
\end{grammarlens}

\subsection*{Your turn}
\begin{itemize}
\raggedright
  \item \emph{Say it:} \emph{kāī}
  \item \emph{Say it:} \emph{kāī}, once more
  \item \emph{Say it:} \emph{pattharoṁ par kāī thī, isliye maiṁ gir gayā.} — There was moss on the stones, so I fell over.
  \item \emph{Recall:} say the Hindi for a bower, then the Hindi for a creeper, then say \emph{kāī} again
\end{itemize}

\begin{tcolorbox}[breakable,colback=teal!4,colframe=teal!35!black,title={Before you move on}]
What does \hi{काई} mean? (\textquotedblleft{}Moss; green scum (on water)\textquotedblright{}.) Say it once more.
\end{tcolorbox}
\section[\texorpdfstring{\hi{कीड़ा} (kīṛā)}{कीड़ा (kīṛā)}]{\hi{कीड़ा} (kīṛā) — an insect; a worm}

\label{lesson:HI-C447-kira}



\begin{quote}
Before the new one: say the Hindi for a creeper, then the Hindi for moss.
\end{quote}

\subsection*{You'll want to know: \hi{कीड़ा}}
\textbf{\hi{कीड़ा}} — \emph{kīṛā} — \textquotedblleft{}an insect; a worm\textquotedblright{}.

\begin{quote}
\hi{पत्ते} \hi{पर} \hi{एक} \hi{हरा} \hi{कीड़ा} \hi{था।}
\end{quote}

\begin{quote}
\emph{patte par ek harā kīṛā thā.}
\end{quote}

\begin{quote}
There was a green insect on the leaf.
\end{quote}

Masculine.

\begin{cousinweb}
\textbf{\hi{कीड़ा}} (\emph{kīṛā}) is a tadbhava of Sanskrit \emph{kīṭa}, \textquotedblleft{}an insect, a worm\textquotedblright{} (CDIAL 3193), which lists Hindi \emph{kīṛā}.
\end{cousinweb}

\begin{grammarlens}[title={What you've built}]
One more word for getting around.
\end{grammarlens}

\subsection*{Your turn}
\begin{itemize}
\raggedright
  \item \emph{Say it:} \emph{kīṛā}
  \item \emph{Say it:} \emph{kīṛā}, once more
  \item \emph{Say it:} \emph{patte par ek harā kīṛā thā.} — There was a green insect on the leaf.
  \item \emph{Recall:} say the Hindi for a creeper, then the Hindi for moss, then say \emph{kīṛā} again
\end{itemize}

\begin{tcolorbox}[breakable,colback=teal!4,colframe=teal!35!black,title={Before you move on}]
What does \hi{कीड़ा} mean? (\textquotedblleft{}An insect; a worm\textquotedblright{}.) Say it once more.
\end{tcolorbox}
\section[\texorpdfstring{\hi{लिपटना} (lipaṭnā)}{लिपटना (lipaṭnā)}]{\hi{लिपटना} (lipaṭnā) — to cling; to wind round}

\label{lesson:HI-C447-lipatna}



\begin{quote}
Before the new one: say the Hindi for moss, then the Hindi for an insect.
\end{quote}

\subsection*{You'll want to know: \hi{लिपटना}}
\textbf{\hi{लिपटना}} — \emph{lipaṭnā} — \textquotedblleft{}to cling; to wind round\textquotedblright{}.

\begin{quote}
\hi{बेल} \hi{पेड़} \hi{से} \hi{लिपट} \hi{गई।}
\end{quote}

\begin{quote}
\emph{bel peṛ se lipaṭ gaī.}
\end{quote}

\begin{quote}
The creeper wound itself round the tree.
\end{quote}

Past (he …): \hi{वह} \hi{बच्चा} \hi{माँ} \hi{से} \hi{लिपट} \hi{गया}, \emph{vah baccā māṁ se lipaṭ gayā}.

\begin{cousinweb}
\textbf{\hi{लिपटना}} (\emph{lipaṭnā}) is filed by CDIAL under Sanskrit \emph{lipyate}, \textquotedblleft{}is smeared\textquotedblright{} (CDIAL 11061), with Hindi \emph{lipaṭnā}, \textquotedblleft{}to stick\textquotedblright{}: to stick close round something.
\end{cousinweb}

\begin{grammarlens}[title={What you've built}]
One more word for getting around.
\end{grammarlens}

\subsection*{Your turn}
\begin{itemize}
\raggedright
  \item \emph{Say it:} \emph{lipaṭnā}
  \item \emph{Say it:} \emph{lipaṭnā}, once more
  \item \emph{Say it:} \emph{bel peṛ se lipaṭ gaī.} — The creeper wound itself round the tree.
  \item \emph{Recall:} say the Hindi for moss, then the Hindi for an insect, then say \emph{lipaṭnā} again
\end{itemize}

\begin{tcolorbox}[breakable,colback=teal!4,colframe=teal!35!black,title={Before you move on}]
What does \hi{लिपटना} mean? (\textquotedblleft{}To cling; to wind round\textquotedblright{}.) Say it once more.
\end{tcolorbox}
